\section{The trivial and quadratic residual characters in the complementary parity}
\label{sz-selmer-section}

The complementary-parity estimate in
Section~\ref{sec:selection-assembly} requires several different
arithmetic arguments. We first prove, without an automorphy input,
the exact vanishing assertion used for its trivial Dirichlet
character. This is the case called $Z$ in the submitted
construction.\cite{alpo-uniform-source} Its proof is shorter than
passing to the full cyclotomic tower: the necessary lifting of
local coboundaries can be done over $\mathbb Q$ itself.

Let $p$ be an odd prime, let $L/\mathbb Q_p$ be finite, and
write $\mathcal O$ for its ring of integers, $\varpi$ for a
uniformizer and $\mathbf k=\mathcal O/\varpi\mathcal O$ for
its residue field. Write $\varepsilon_p:G_{\mathbb Q}\to
\mathbb Z_p^\times$ for the cyclotomic character. For $e\ge1$
put
\[
 A_e=(L/\mathcal O)(-e),\qquad
       g\cdot a=\varepsilon_p(g)^{-e}a,\qquad
 B_e=A_e[\varpi]\simeq\mathbf k(\bar\varepsilon_p^{-e}).
                                                               \tag{SZ1}
\]
These are discrete modules. Fix decomposition groups $D_q$
and inertia groups $I_q$ at the rational primes $q$. Define
\[
 \mathcal T_p(e)=\ker\left[
 H^1(G_{\mathbb Q},A_e)\longrightarrow
 H^1(D_p,A_e)\oplus\prod_{q\ne p}H^1(I_q,A_e)\right].   \tag{SZ2}
\]
This is the strict Selmer group with local condition zero
at $p$ and unramified conditions elsewhere. Adding zero
local conditions at further primes only replaces it by a
subgroup.

\subsection{Lifting a torsion class with its local condition}

We use the following standard form of unramified class-field
reciprocity: continuous homomorphisms from $G_F$ to a finite
abelian group of odd order that vanish on all finite inertia
groups factor through the ordinary ideal class group of $F$.
At a real place the inertia group has order at most two and
therefore has zero image. This is the Hilbert class field
form of the global reciprocity theorem.\cite[Chapter V,
\S3]{milne-cft}

\begin{lemma}\label{sz-torsion-and-local}
Every element of $H^1(G_{\mathbb Q},A_e)$ is killed by a
power of $\varpi$. The image of the cyclotomic character
on $D_p$ equals its image on $G_{\mathbb Q}$, namely
$\mathbb Z_p^\times$. Consequently
\[
                         A_e^{G_{\mathbb Q}}=A_e^{D_p}.
                                                               \tag{SZ3}
\]
Every class in $\mathcal T_p(e)[\varpi]$
is the image of a class $c\in H^1(G_{\mathbb Q},B_e)$
whose restrictions to $D_p$ and all $I_q$, $q\ne p$,
are zero.
\end{lemma}
\begin{proof}
A continuous cocycle from a profinite group to the discrete
module $A_e$ has finite image. Choose a common power of
$\varpi$ killing those finitely many values; it kills the
cohomology class as well.

For $r\ge1$, the polynomial $\Phi_{p^r}(1+X)$ is Eisenstein
at $p$. Its constant term is $p$, and its reduction modulo
$p$ is $X^{p^{r-1}(p-1)}$. Hence
\[
 [\mathbb Q_p(\mu_{p^r}):\mathbb Q_p]=p^{r-1}(p-1).
\]
Its Galois group therefore gives the full group
$(\mathbb Z/p^r\mathbb Z)^\times$ under the cyclotomic
character. Taking inverse limits gives the local assertion;
the global image contains the local one. Formula (SZ3)
follows because the action on $A_e$ factors through that
character.

The exact sequence
\[
 0\longrightarrow B_e\longrightarrow A_e
          \xrightarrow{\,\varpi\,}A_e\longrightarrow0    \tag{SZ4}
\]
shows that a class $x\in H^1(G_{\mathbb Q},A_e)[\varpi]$
lifts to $c\in H^1(G_{\mathbb Q},B_e)$. If $x\in
\mathcal T_p(e)$, the restriction $c|_{D_p}$ belongs to
the image of the connecting map
\[
 A_e^{D_p}/\varpi A_e^{D_p}\longrightarrow H^1(D_p,B_e).
\]
By (SZ3) choose an element of
$A_e^{G_{\mathbb Q}}/\varpi A_e^{G_{\mathbb Q}}$ with
this same image. Subtract its global connecting class from
$c$. Its image in $H^1(G_{\mathbb Q},A_e)$ is unchanged,
and the new $c$ has zero restriction to $D_p$.

For $q\ne p$ the inertia group acts trivially on $A_e$,
since the cyclotomic character is unramified there.
Multiplication by $\varpi$ is surjective on $A_e$, so
(SZ4) shows that
\[
                     H^1(I_q,B_e)\longrightarrow H^1(I_q,A_e)
\]
is injective. The image of $c$ on the right is zero by
(SZ2), so its restriction on the left is zero as well.
The stated lifting conclusion follows.
\end{proof}

\begin{proposition}\label{sz-residual-class}
Let $F_e$ be the fixed field in $\mathbb Q(\mu_p)$ of the
kernel of $\bar\varepsilon_p^{-e}$. There is an injection
\[
       \mathcal T_p(e)[\varpi]\lhook\joinrel\longrightarrow
                  \operatorname{Hom}(\operatorname{Cl}(F_e),B_e),
                                                               \tag{SZ5}
\]
where the target denotes homomorphisms of additive abelian
groups, and $B_e$ is trivial as a $G_{F_e}$-module.
In particular $p\nmid h_{F_e}$ implies $\mathcal T_p(e)=0$.
Moreover
\[
 \dim_{\mathbf k}\mathcal T_p(e)[\varpi]
       \le\dim_{\mathbb F_p}\operatorname{Cl}(F_e)/p.
                                                               \tag{SZ6}
\]
\end{proposition}
\begin{proof}
The final part of Lemma~\ref{sz-torsion-and-local} constructs
a lift with zero restriction at $D_p$ for every $e$.
It is unique: two lifts differ by a global connecting class
from $A_e^{G_{\mathbb Q}}/\varpi A_e^{G_{\mathbb Q}}$.
The local connecting map is injective on
$A_e^{D_p}/\varpi A_e^{D_p}$ by exactness, and (SZ3)
identifies its source with the global one. A connecting
class whose local restriction is zero is therefore zero.
Thus this lift depends $\mathcal O$-linearly on the original
class.

Restrict to $G_{F_e}$. This is injective on
$H^1(G_{\mathbb Q},B_e)$: the inflation kernel is
$H^1(\operatorname{Gal}(F_e/\mathbb Q),B_e)$, which
vanishes by averaging because its order divides $p-1$.
On $G_{F_e}$ the cocycle is a homomorphism. It is
unramified everywhere by the local conditions on its lift,
so class-field reciprocity gives (SZ5). A finite abelian
group with $p$-rank $r$ has a space of homomorphisms to
$\mathbf k$ of $\mathbf k$-dimension $r$, proving (SZ6).
If $p\nmid h_{F_e}$, (SZ5) is zero; the torsion argument
used below then annihilates the full group.
\end{proof}

\begin{theorem}\label{sz-trivial}
If $p-1\mid e$, then $\mathcal T_p(e)=0$.
\end{theorem}
\begin{proof}
The residual action on $B_e$ is trivial. A class furnished
by Lemma~\ref{sz-torsion-and-local} is therefore a continuous
homomorphism $G_{\mathbb Q}\to(B_e,+)$ which is unramified
at every finite prime. Its target is a finite $p$-group.
Class-field reciprocity makes it factor through
$\operatorname{Cl}(\mathbb Q)=0$. Thus
$\mathcal T_p(e)[\varpi]=0$. A nonzero element of a
$\varpi$-primary torsion module has a nonzero multiple
killed by $\varpi$; the first assertion of the lemma
therefore proves $\mathcal T_p(e)=0$.
\end{proof}

\subsection{The quadratic residual character}

\begin{lemma}\label{sz-small-real-class}
If $p\equiv1\pmod4$ is prime and $F=\mathbb Q(\sqrt p)$,
then $h_F<p$.
\end{lemma}
\begin{proof}
The discriminant is $p$. The small-ideal and ideal-count
bounds in Section~\ref{sl-section}, with $[F:\mathbb Q]=2$,
give
\[
                         h_F\le\sqrt p\,(1+\log\sqrt p).
\]
For $u>1$, integrating $1/t<1$ from $1$ to $u$ gives
$\log u<u-1$. Taking $u=\sqrt p$ makes the right side
strictly less than $p$.
\end{proof}

\begin{theorem}\label{sz-quadratic}
Suppose $e$ is even, $p-1\mid2e$, and $p-1\nmid e$.
Then $\mathcal T_p(e)=0$.
\end{theorem}
\begin{proof}
Write $p-1=2a$. The divisibilities say $e=ab$ with $b$
odd. Since $e$ is even, $a$ is even and $p\equiv1\pmod4$;
in particular $p\ge5$. The residual character on $B_e$
is the quadratic character
$\bar\varepsilon_p^{(p-1)/2}$. Its fixed field is
$F=\mathbb Q(\sqrt p)$. One explicit identification is the
quadratic Gauss sum
\[
 \tau=\sum_{u\bmod p}(u/p)\zeta_p^u,\qquad
 \sigma_v\tau=(v/p)\tau,\qquad \tau^2=p.
\]
The last identity follows from
$\tau\bar\tau=p$ and $(-1/p)=1$, as in the quadratic
field calculation of Section~\ref{sec:selection-assembly}.

There is an element of $D_p$ on which this residual
character is $-1$. Its action minus one on $A_e$ is
multiplication by a unit of $\mathcal O$. Hence
$A_e^{D_p}=0$. Given $x\in\mathcal T_p(e)[\varpi]$,
choose its lift $c\in H^1(G_{\mathbb Q},B_e)$ using
(SZ4). The local connecting term at $p$ is now zero,
so $c|_{D_p}=0$. As in the preceding lemma, its
restrictions to every $I_q$, $q\ne p$, are zero.

Restriction to $G_F$ makes $B_e$ a trivial module.
Thus $c|_{G_F}$ is an everywhere unramified homomorphism
from $G_F$ to the additive group of $\mathbf k$.
At primes above $p$ it is zero even on the decomposition
groups; at all other finite primes inertia is already
killed. Class-field reciprocity makes it factor through
$\operatorname{Cl}(F)$. Lemma~\ref{sz-small-real-class}
shows $p\nmid h_F$, so this homomorphism is zero.

The restriction map
$H^1(G_{\mathbb Q},B_e)\to H^1(G_F,B_e)$ is injective.
Indeed its kernel is the image of
$H^1(\operatorname{Gal}(F/\mathbb Q),B_e)$ by the
inflation--restriction sequence, and this latter group
vanishes by averaging over its group of order two; two
is invertible on $B_e$. Thus $c=0$ and $x=0$.
The torsion argument in Theorem~\ref{sz-trivial} finishes
the proof.
\end{proof}

\begin{corollary}\label{sz-exact-case}
For every prime $p\ge5$, every even $e\ge2$ such that
$p-1\mid2e$, and every finite extension $L/\mathbb Q_p$,
one has $\mathcal T_p(e)=0$. The same is true if additional
zero local conditions are imposed at any finite set of
primes. In particular the length for the trivial Dirichlet
character in case $Z$ of the submitted complementary-parity
argument is exactly zero.
\end{corollary}
\begin{proof}
Apply Theorem~\ref{sz-trivial} when $p-1\mid e$ and
Theorem~\ref{sz-quadratic} otherwise. An additional local
condition cuts out a subgroup, so it preserves vanishing.
The source's coefficient action for the trivial character
is exactly $\varepsilon_p^{-e}$, and its level-$N$ group
adds zero conditions at the primes dividing $N$.
Thus its group is a subgroup of (SZ2). Its parity condition
at the trivial character is precisely that $e$ be even.
\end{proof}

This proves one complete arithmetic input, including
arbitrarily ramified coefficient extensions and all the
relevant twists. It does not identify the groups of
nontrivial Dirichlet characters with these trivial-character
groups. The complementary-parity estimate for the other
characters still requires the remaining counting and
restriction arguments, with constants independent of the
selected auxiliary primes.


